\section{Future Work}
\label{sec:future-work}

The next step is not to claim operational readiness, but to broaden external
validity while preserving the safety-control interpretation developed here.
Future studies should test richer station layouts, sensor assumptions, hazard
models, and calibrated population mixes. They should also compare Chiyoda's
stylized trajectories against evacuation drills, incident records, controlled
VR experiments, and pedestrian trajectory datasets before making predictive
claims about specific stations or events.

Adaptive message generation is a second extension, including learned or
language-model-mediated guidance. Section~\ref{subsec:llm-guidance-evaluation}
shows that generated guidance can be replayable and safety-efficient under a
sparse budget, but it also shows why messages should not be evaluated by
fluency alone. Future generated-language systems should remain replayable,
validated against the simulated state, prevented from inventing unavailable
exits or stale hazard claims, and measured by downstream exposure, queueing,
exit balance, and evacuation outcomes.
